\section{Conclusion}
The relevant dimension of joint renewal design is not necessarily the number of observed histories. Common catalog eligibility permits exact pooling of arbitrary participation caps and quadratic service costs. Terminal-first allocation isolates the book dependence of heterogeneous service in a last-eligible-level correction. Paying that correction on the ordered path gives a polynomial exact joint algorithm for one eligibility class and an exact group-box price oracle for several classes.

The resulting certificates retain the original command budget, opening charges, aggregate promise, and individual constraints. Their adaptive accuracy dependence is controlled by eligibility complexity rather than history count or dispersion bins. Exact single-class computations and checked multigroup intervals are reported separately; the latter include the unsuccessful accuracy requests. The preserved harmonic, continuous, institutional, Monge, compression, and inexact-support results describe additional regimes rather than substitute for the new exact joint theorem. The study establishes algorithmic and certificate behavior on specified synthetic inputs, not a calibrated field effect.

\paragraph{Code and data availability.}
The accompanying repository contains the immutable review and scientific-parent identifiers, the pre-execution protocol, all rational instances and policy certificates, independent checkers, numerical mixed-integer records, source files, and build instructions. The computational record identifies every case, including interruptions, and separates historical, local-validation, and publication-run timing environments. A single certificate can be checked without rerunning any optimizer. The terminal alphabet count is not a claim about total implementation memory: target tables, encoder logic, read-only constants, probability precision, and certificate storage remain distinct resources.
